\documentclass[12pt]{article}

\usepackage[a4paper, total={6in, 8in}]{geometry}
\usepackage{textcomp}

\begin{document}
\setlength{\parindent}{0pt}

\def \t {\textrightarrow}
\def \v {\vspace{1em}}
\def \bi {\begin{itemize}}
\def \ei {\end{itemize}}
\def \s[#1] {\section*{#1}}
\def \ss[#1] {\subsection*{#1}}
\def \sss[#1] {\subsubsection*{#1}}

\s[Recuperare 24 febbraio]
\s[Proteine]
Cos'è una proteina? \t è un polimero, il cui monomero è l'amminoacido che ha gruppo carbossilico (COOH) e ammidico \\
Si ha stereocentro di simmetria chirale in corrispondenza del carbonio $\alpha$ \\
Ci sono legami peptidici, che coinvolgono un gruppo carbossilico e un gruppo ammidico di due diversi amminoacidi \\
Le proteine si differenziano per la combinazione di amminoacidi, che sono 20 \\
Gli enzimi sono proteine, che fanno da catalizzatori \\
Sapere le varie funzioni: catalitica, di trasporto, strutturale, difesa, segnalazione, etc. \\
Proteina semplice è fatta solo di amminoacidi \\
Proteina coniugata = ha gruppi prostetici, come carboidrati o altre molecole organiche, o ioni metallici \\
Proteina fibrosa ha funzione strutturale: ha struttura lineare assembabile, mentre una proteina globulare no: non ha compatezza \\
Proteine fibrose sono insolubili, ma possono esserlo a condizioni particolari (alte temperature etc.) \t es. cheratina \\
Mioglobina è globulare, porta ossigeno nei muscoli e quindi deve essere solubile

\v

Polarita determina la solubilita in acqua \t pero cambia sempre in base alla lunghezza della catena idrofobica, che puo contrastare il gruppo elettronegativo \\
Zwitterione \t gruppo NH2 e COOH \t l'amminoacidi sono presenti nella stessa molecola, quindi ha comportamento acido/base, che dipenderà dal ph dell'ambiente (se ho ambiente acido, molecola si comportera come base ??) \\
In questo caso i valori del ph sono abbastanza neutri, quindi ho uno ione dipolare \\
\textbf{Definizione di acido: è un composto che in soluzione acquosa rilascia ioni H+} \t ambiente acido è ambiente con tanti ioni H+ \\
Quindi ioni H+ si legeranno + probabilmente alla parte basica con lo zwitterione

\v

Punto isoelettrico \t che è spesso intorno alla neutralita (ph 7) \\
La prolina è apolare, perche i due gruppi + e - si neutralizzano \t se la catena laterale è neutra, allora è apolare (alifatico non è aromatico) \\
Con catena aromatica invece è apolare perchè è molto stabile \\
Gruppi R carichi negativamente allora sono acidi, R positivi basici

\v

Perche parliamo di amminoacidi essenziali? \t con questi si indicano gli amminoacidi che servono al corpo umano ma non li puo sintetizzare \t e sono 8 \\
Con la dieta dobbiamo introdurre le vitamine e gli amminoacidi essenziali \\
La cisteina altera la proteina \t e la reazione però è reversibile
\end{document}
